% PDF-only preamble for the Kates book, pulled in by format.pdf's
% include-in-header in _quarto.yml. The web theme never sees it; the colours
% are the web palette's (kates-light.scss), so both editions read as one book.

% ── Colours ─────────────────────────────────────────────────────────────────
\definecolor{katesblue}{HTML}{1D63C7}
\definecolor{katesnavy}{HTML}{131D42}
\definecolor{katesmuted}{HTML}{5B6683}
\definecolor{katesoutputbg}{HTML}{F6F8FC}
\definecolor{katesoutputfg}{HTML}{26324D}
\definecolor{katesoutputrule}{HTML}{9FB3D6}

% ── Fonts ───────────────────────────────────────────────────────────────────
% The web's type: Inter for text and headings, JetBrains Mono for code. Each
% face is used only when it is installed, so a render on a machine without
% them still works, in the next font on the list (Latin Modern if none is).
% Characters a face lacks (check marks, triangles, warning signs, the Lab
% mockup's emoji) come from a fallback list instead of vanishing. The pdf job
% in .github/workflows/book.yml installs Inter, JetBrains Mono, the Noto
% symbol fonts and Symbola, and fails on any "Missing character" in the log.
\ifLuaTeX
\makeatletter
\newcommand*\kates@pick[2]{% #1: macro to set, #2: candidate faces, in order
  \let#1\relax
  \@for\kates@face:=#2\do{%
    \ifx#1\relax\IfFontExistsTF{\kates@face}{\let#1\kates@face}{}\fi}}
\kates@pick\kates@text{Inter,Helvetica Neue,Arial,DejaVu Sans}
\kates@pick\kates@code{JetBrains Mono,DejaVu Sans Mono,Menlo}
% Symbols, then emoji drawn in black and white. The last two are macOS system
% fonts, so a local render loses fewer characters.
\def\kates@fallbacks{}
\@for\kates@face:={Noto Sans Symbols 2,Noto Sans Symbols,Noto Sans Math,%
    Symbola,Apple Symbols,Arial Unicode MS}\do{%
  \IfFontExistsTF{\kates@face}{%
    \edef\kates@fallbacks{\kates@fallbacks"\kates@face:",}}{}}
\ifx\kates@fallbacks\@empty
  \def\kates@features{}
\else
  \directlua{luaotfload.add_fallback("katesfallback", {\kates@fallbacks})}
  \def\kates@features{RawFeature={fallback=katesfallback}}
\fi
\ifx\kates@text\relax\else
  \expandafter\setmainfont\expandafter{\kates@text}[\kates@features]
  \expandafter\setsansfont\expandafter{\kates@text}[\kates@features]
\fi
\ifx\kates@code\relax\else
  \expandafter\setmonofont\expandafter{\kates@code}[\kates@features]
\fi
\makeatother
\fi

% ── Page ────────────────────────────────────────────────────────────────────
% Let a page end short rather than stretch the space between paragraphs to
% fill it: with scrbook's default \flushbottom, a diagram that moves to the
% next page leaves its paragraphs pulled far apart.
\raggedbottom

% Paragraphs are set apart by space, as before, but without KOMA's rule that
% a paragraph's last line keep 1em free: a diagram at full line width is a
% paragraph of one line, and the rule pushed each one 1em into the margin.
\KOMAoptions{parskip=half-}

% Section and figure numbers without a trailing period ("2.1", "Figure 2.1:"),
% as on the web.
\KOMAoptions{numbers=noendperiod}

% Let figures share a page with text rather than sit alone on a float page
% with a gap above and below.
\renewcommand{\topfraction}{0.85}
\renewcommand{\bottomfraction}{0.7}
\renewcommand{\textfraction}{0.1}
\renewcommand{\floatpagefraction}{0.75}

% Headings, the title and links in the brand blue; running heads and page
% numbers recede.
\addtokomafont{disposition}{\color{katesblue}}
\addtokomafont{title}{\color{katesblue}}
\addtokomafont{subtitle}{\color{katesnavy}}
\addtokomafont{pagehead}{\normalfont\small\color{katesmuted}}
\addtokomafont{pagenumber}{\normalfont\small\color{katesmuted}}

% Part titles already say "Part I — Foundations", so drop KOMA's own
% "Part I." label above them, and its "I" in the table of contents.
\renewcommand*{\partformat}{}
\renewcommand*{\addparttocentry}[2]{\addtocentrydefault{part}{}{#2}}

% Room in the table of contents for three-digit page numbers and for section
% numbers such as 26.10: Inter's digits are wider than the defaults allow for.
\makeatletter
\renewcommand*{\@pnumwidth}{2.2em}
\makeatother
\RedeclareSectionCommand[tocnumwidth=2.9em]{section}

% ── Graphics ────────────────────────────────────────────────────────────────
% Mermaid diagrams arrive at their natural size, often wider than the page.
% Shrink any graphic to the line width and to 80% of the text height, keeping
% its proportions; smaller ones stay as they are. adjustbox applies its max
% keys after the width and height Pandoc writes, whatever their order.
% \kates@includegraphics is the unclamped original, for the cover.
\usepackage[export]{adjustbox}
\makeatletter
\NewCommandCopy\kates@includegraphics\includegraphics
\RenewDocumentCommand\includegraphics{s O{} m}{%
  \IfBooleanTF{#1}%
    {\kates@includegraphics*[#2,max width=\linewidth,max totalheight=0.8\textheight]{#3}}%
    {\kates@includegraphics[#2,max width=\linewidth,max totalheight=0.8\textheight]{#3}}}
\makeatother

% ── Tables ──────────────────────────────────────────────────────────────────
% Tables one step smaller than the text, so their columns, and the code in
% them, get more room.
\AddToHook{env/longtable/begin}{\small}

% ── Code and output ─────────────────────────────────────────────────────────
% Commands and files on a light tint with a brand-blue rule. Pandoc defines
% Shaded before this header is read; this replaces it. The code is two steps
% smaller than the text, so about 76 characters of JetBrains Mono fit on a
% line.
\usepackage{fancyvrb}
\usepackage{tcolorbox}
\tcbuselibrary{skins,breakable}
\definecolor{katescodebg}{HTML}{EEF2F9}
\makeatletter
\@ifundefined{Shaded}{}{%
  \renewenvironment{Shaded}{%
    \begin{tcolorbox}[enhanced, breakable, sharp corners, frame hidden,
      boxrule=0pt, colback=katescodebg,
      borderline west={2.5pt}{0pt}{katesblue},
      boxsep=0pt, left=8pt, right=6pt, top=5pt, bottom=5pt,
      before skip=10pt, after skip=10pt]}%
    {\end{tcolorbox}}}
\@ifundefined{Highlighting}{}{%
  \RecustomVerbatimEnvironment{Highlighting}{Verbatim}%
    {commandchars=\\\{\},fontsize=\footnotesize}}
\makeatother

% `text` fences (terminal output and ASCII screens), as theme/pdf-output.lua
% sets them: a lighter panel with a thin grey-blue rule, the web's output
% colours, so what you type and what you see never look alike. The filter
% adds [output label=Output] when the fence is output; the label sits in the
% panel's top margin. Lines sit exactly a baseline apart, however tall their
% glyphs, so the box-drawing lines of a mockup meet instead of breaking.
\tcbset{output label/.style={top=18pt,
  overlay unbroken and first={\node[anchor=north west, inner sep=0pt,
    font=\sffamily\scriptsize\bfseries, text=katesmuted]
    at ([xshift=8pt, yshift=-5pt]frame.north west) {\MakeUppercase{#1}};}}}
\newtcolorbox{katesoutput}[1][]{%
  enhanced, breakable, sharp corners, frame hidden,
  boxrule=0pt, colback=katesoutputbg, coltext=katesoutputfg,
  borderline west={1.5pt}{0pt}{katesoutputrule},
  boxsep=0pt, left=8pt, right=6pt, top=5pt, bottom=5pt,
  before skip=10pt, after skip=10pt,
  fontupper=\footnotesize,
  before upper={\lineskiplimit=-\maxdimen},
  #1}

% ── Title page ──────────────────────────────────────────────────────────────
% When theme/pdf-cover.lua has drawn assets/cover.svg into a PDF, the book
% opens on it, full bleed, before the usual title page.
\makeatletter
\let\kates@maketitle\maketitle
\renewcommand{\maketitle}{%
  \ifdefined\katescoverfile
    \begingroup
    \thispagestyle{empty}%
    \AddToHookNext{shipout/background}{%
      \put(0,-\paperheight){%
        \kates@includegraphics[width=\paperwidth,height=\paperheight]{\katescoverfile}}}%
    \null\clearpage
    \endgroup
    \cleardoublepage
  \fi
  \kates@maketitle}
\makeatother
